\chapter{基于蕴含锥的双曲层次图神经网络}

本章采用第三章提出的显式层次建模范式，通过蕴含锥机制直接约束路网层次关系，解决第四章隐式建模中层次包含关系缺乏几何约束的不足。

\section{研究动机与问题分析}

第四章的 HRGNN 证明了双曲空间对路网层次建模的有效性，但其隐式范式存在两个根本性不足。其一，层次关系的表达缺乏方向性。双曲距离可以反映两个嵌入点的层次接近程度，但无法区分"父-子"关系和"兄弟"关系——一个功能区嵌入与其下属路段嵌入之间的双曲距离，可能与两个同层级路段嵌入之间的距离相当，导致层次包含关系的语义模糊。其二，HRGNN 仅在路段层级上工作，不构建局部区域和功能区的显式表示，无法支持需要多层次信息交互的通用表示学习任务。

这些不足指向一个明确的改进方向：在保持双曲空间几何优势的前提下，引入专门的几何工具对层次包含关系施加显式约束，并构建多层次的结构化表示。

HRNR\cite{wu2020hrnr} 构建了路段、结构区域和功能区三层次的神经架构，通过谱聚类将路段划分为结构区域，通过轨迹数据的协同过滤学习功能区的软性划分。HRNR 的成功表明对路网层次结构的显式建模带来了可观的性能收益，但其在欧氏空间中运行，层次关系通过赋值矩阵以软性方式表达，缺乏几何意义上的层次约束。路段嵌入与其所属结构区域的嵌入之间，仅通过矩阵乘法和注意力机制关联，而非利用嵌入空间的几何结构显式约束"包含"关系。

HiFiNet 通过引入虚拟层次节点来增强 GNN 的层次感知能力，避免了复杂的聚类预处理，但层次关系的表达仍然是隐式的。

路网中存在"通用-具体"层次关系。功能区（Region）覆盖大片地理空间，对应"最通用"概念，数量少（$|\mathcal{R}| \approx 30$）；局部区域（Locality）由地理邻近路段聚合而成，对应"中等通用"概念（$|\mathcal{C}| \approx 300$）；路段（Segment）为个体路段，对应"最具体"概念，数量最多。这种"Region $\supseteq$ Locality $\supseteq$ Segment"的包含关系。

近年来，Zhou 等\cite{zhou2024garner} 通过引入地理学第三定律改进路网表示的空间语义编码，展示了地理规律先验对路网嵌入质量的提升潜力；然而该方法仍在欧氏空间中运行，无法从几何上利用路网层次结构。Sun 等\cite{sun2024lsenet} 在图聚类任务中利用 Lorentz 结构熵指导层次划分，为双曲空间中多层次图组织提供了实证支撑。这些工作进一步证实了路网层次显式建模的重要价值，以及双曲几何在此类任务中的适用性。

蕴含锥在双曲空间中具有三个适用于路网的关键性质。靠近原点的点（功能区嵌入）具有宽蕴含锥，几何上"包含"更多子概念，与功能区覆盖多个局部区域的语义一致。远离原点的点（路段嵌入）具有窄蕴含锥，代表高度特定的概念。蕴含锥损失通过明确的角度约束强制层次包含关系，实现了从软赋值矩阵到硬几何约束的质变。

如第三章分析，Lorentz 模型的指数映射和对数映射具有封闭形式，数值稳定性优于庞加莱球模型。本章的模型需要在三层层次结构中精确组织嵌入的几何关系，路段嵌入需要位于远离原点处（对应庞加莱球的边界附近），这对数值稳定性提出了更高要求。此外，Lorentz 模型中到原点的距离天然编码层次深度——功能区嵌入于靠近原点处，路段嵌入于远离原点处——这种几何安排在优化目标包含层次信号时自然出现，为蕴含锥约束提供了良好的初始条件。

\section{模型架构设计}

HHRoad 的目标是为路网中每条路段学习一个双曲嵌入向量，使其同时编码路段自身的属性特征、拓扑邻接关系以及所属的多层次结构信息。本节详细阐述 HHRoad 的完整架构设计，模型的数据流从路段原始属性出发，经过双曲嵌入层、层次划分与表示聚合、自顶向下双曲消息传递三个核心阶段，最终产出融合了三层层次语义的路段双曲嵌入。图~\ref{fig:HHRoad_arch} 给出了模型的整体架构示意。

\begin{figure}[H]
\centering
\includegraphics[width=\textwidth]{figures/HHRoad模型架构.png}
\caption{HHRoad 模型架构总览}
\label{fig:HHRoad_arch}
\end{figure}

\subsection{整体框架}

HHRoad 由五个核心组件构成：

\textbf{（1）双曲嵌入层}将每条路段的多维欧氏属性向量通过可学习的线性变换和指数映射投影至 Lorentz 双曲面 $\mathbb{L}^n$，产出路段级初始双曲嵌入。

\textbf{（2）层次划分与表示聚合}基于谱聚类和轨迹共现模式分别构建路段到局部区域、局部区域到功能区的软赋值矩阵，并在 Lorentz 切空间中自底向上逐级聚合，产出局部区域嵌入和功能区嵌入。至此，路网的三层层次表示体系（路段 $\to$ 局部区域 $\to$ 功能区）构建完成。

\textbf{（3）自顶向下双曲消息传递}遵循 R2R $\to$ R2C $\to$ C2C $\to$ C2S $\to$ S2S 的五步传播方案，在 Lorentz 空间中依次进行功能区间消息传递、功能区向局部区域分发、局部区域间消息传递、局部区域向路段分发和路段间消息传递，使全局层次语义沿层次结构高效下行至每条路段。

\textbf{（4）组合蕴含学习}通过蕴含锥机制在双曲空间中对三类层次包含关系施加显式几何约束，确保功能区嵌入的蕴含锥包含其下属局部区域、局部区域的蕴含锥包含其下属路段、拓扑相邻路段具有互兼容的蕴含锥结构。

\textbf{（5）层次对比学习}在路段层内以拓扑相邻路段为正对、跨层以路段与其父级局部区域为正对，分别通过 InfoNCE 损失强化嵌入空间中的局部语义聚类和跨层次关联。

总训练损失由结构重构损失 $\mathcal{L}_{\text{struct}}$、轨迹重构损失 $\mathcal{L}_{\text{traj}}$、蕴含损失 $\mathcal{L}_{\text{CE}}$ 和对比损失 $\mathcal{L}_{\text{CC}}$ 的加权和构成（详见第~\ref{sec:loss} 节）。以下依次阐述各组件的具体设计。

\subsection{双曲嵌入层}

双曲嵌入层的任务是将路段的欧氏属性向量投影为 Lorentz 双曲面上的初始嵌入，为后续所有双曲空间操作提供起点。

给定路段 $v_i$ 的原始属性集合——节点编号、道路类型、路段长度和车道数，首先通过各自的嵌入查找表获取独立的特征向量，然后拼接为统一的欧氏属性向量：
\begin{equation}
    \mathbf{x}_i = [\mathrm{Emb}_{\text{node}}(v_i) \| \mathrm{Emb}_{\text{type}}(v_i) \| \mathrm{Emb}_{\text{length}}(v_i) \| \mathrm{Emb}_{\text{lane}}(v_i)] \in \mathbb{R}^{d_{\text{in}}}
\end{equation}
其中 $\|$ 表示向量拼接，$d_{\text{in}}$ 为拼接后的总维度。随后通过可学习线性变换 $f_\theta: \mathbb{R}^{d_{\text{in}}} \to \mathbb{R}^n$ 将其映射为 $n$ 维向量，再经由 Lorentz 投影将空间分量补全为满足双曲约束的 $(n+1)$ 维向量：
\begin{equation}
    \mathbf{z}_i = f_\theta(\mathbf{x}_i) \in \mathbb{R}^n, \quad \mathbf{h}_i = \left(\sqrt{1/\kappa + \|\mathbf{z}_i\|^2},\ \mathbf{z}_i\right)^\top \in \mathbb{L}^n
\end{equation}
其中时间分量 $h_{i,0} = \sqrt{1/\kappa + \|\mathbf{z}_i\|^2}$ 确保 $\mathbf{h}_i$ 满足 Lorentz 约束 $\langle \mathbf{h}_i, \mathbf{h}_i \rangle_{\mathbb{L}} = -1/\kappa$。经过此步骤，全部 $|\mathcal{S}|$ 条路段的初始双曲嵌入矩阵 $\mathbf{H}^{(0)} = [\mathbf{h}_1, \ldots, \mathbf{h}_{|\mathcal{S}|}]^\top \in \mathbb{R}^{|\mathcal{S}| \times (n+1)}$ 即构建完成。

\subsection{层次划分与表示聚合}

在路段初始嵌入的基础上，HHRoad 构建"路段 $\to$ 局部区域 $\to$ 功能区"的三层层次表示体系。两个关键要素是层次划分方案和双曲空间中的聚合机制。

\subsubsection{层次划分方案}

局部区域（Locality）的划分基于路网拓扑结构的谱聚类\cite{von2007tutorial}。对路网邻接矩阵的归一化图拉普拉斯矩阵进行特征分解，取前 $|\mathcal{C}|$ 个特征向量（$|\mathcal{C}|=300$），通过 k-means 将路段划分为 $|\mathcal{C}|$ 个结构性子区域，初始化路段到局部区域的软赋值矩阵 $\mathbf{A}_{SC} \in \mathbb{R}^{|\mathcal{S}| \times |\mathcal{C}|}$，行归一化使每行之和为 1，即每条路段对所属局部区域的隶属概率分布。

功能区（Region）的划分基于轨迹数据的协同过滤学习。通过挖掘轨迹中路段的共现模式，学习局部区域到功能区的软赋值矩阵 $\mathbf{A}_{CR} \in \mathbb{R}^{|\mathcal{C}| \times |\mathcal{R}|}$（$|\mathcal{R}|=30$）。两个赋值矩阵均在训练过程中以梯度更新方式进一步微调，以在端到端学习中逐步优化层次划分的质量。

\subsubsection{Lorentz 切空间聚合}

由于 Lorentz 双曲面是弯曲流形，直接在流形上做加权平均不满足封闭性（加权组合不一定仍在 $\mathbb{L}^n$ 上）。HHRoad 采用"对数映射 $\to$ 切空间加权平均 $\to$ 指数映射"的三步聚合策略：首先将子级嵌入通过对数映射投影到 Lorentz 原点 $\mathbf{o} = (\sqrt{1/\kappa}, 0, \ldots, 0)^\top$ 处的切空间 $\mathcal{T}_\mathbf{o}\mathbb{L}^n$，在该切空间（作为欧氏空间）中执行加权平均，最后通过指数映射将结果投影回流形。

局部区域嵌入 $\mathbf{h}_{c_j}$ 由路段嵌入聚合得到：
\begin{equation}
    \mathbf{h}_{c_j}^{\text{tan}} = \sum_{i=1}^{|\mathcal{S}|} [\mathbf{A}_{SC}]_{ij} \cdot \log_{\mathbf{o}}^\kappa(\mathbf{h}_i), \quad \mathbf{h}_{c_j} = \exp_{\mathbf{o}}^\kappa(\mathbf{h}_{c_j}^{\text{tan}})
    \label{eq:locality_agg}
\end{equation}
其中 $\log_{\mathbf{o}}^\kappa(\cdot)$ 为 Lorentz 原点处的对数映射，$\exp_{\mathbf{o}}^\kappa(\cdot)$ 为原点处的指数映射（具体公式见第三章 3.2.3 节），$[\mathbf{A}_{SC}]_{ij}$ 为路段 $v_i$ 对局部区域 $c_j$ 的赋值权重。

功能区嵌入 $\mathbf{h}_{r_k}$ 类似地从局部区域嵌入聚合得到：
\begin{equation}
    \mathbf{h}_{r_k}^{\text{tan}} = \sum_{j=1}^{|\mathcal{C}|} [\mathbf{A}_{CR}]_{jk} \cdot \log_{\mathbf{o}}^\kappa(\mathbf{h}_{c_j}), \quad \mathbf{h}_{r_k} = \exp_{\mathbf{o}}^\kappa(\mathbf{h}_{r_k}^{\text{tan}})
    \label{eq:region_agg}
\end{equation}

这一自底向上的两级聚合完成后，模型获得了三层层次嵌入：路段嵌入 $\{\mathbf{h}_i\}_{i=1}^{|\mathcal{S}|}$、局部区域嵌入 $\{\mathbf{h}_{c_j}\}_{j=1}^{|\mathcal{C}|}$ 和功能区嵌入 $\{\mathbf{h}_{r_k}\}_{k=1}^{|\mathcal{R}|}$，三者均位于同一 Lorentz 空间 $\mathbb{L}^n$ 中。由于聚合过程涉及多次加权平均，高层嵌入的空间范数（$\|\tilde{\mathbf{h}}\|$，即除去时间分量后的欧氏范数）倾向于小于底层嵌入的范数，即功能区嵌入自然靠近 Lorentz 原点、路段嵌入远离原点。这一范数递减特性与蕴含锥的层次语义一致——靠近原点的点具有更宽的蕴含锥（代表更一般的概念），为后续蕴含约束提供了良好的几何初始条件。

\subsection{自顶向下双曲消息传递}
\label{sec:message_passing}

层次表示聚合构建了三层嵌入的初始版本，但各层次嵌入仅编码了自身及其下属的信息。自顶向下消息传递的目标是使全局层次语义沿"功能区 $\to$ 局部区域 $\to$ 路段"的路径逐级下行，使每条路段的最终嵌入不仅包含局部拓扑信息，还融合了其所属功能区和局部区域的宏观语义。

\subsubsection{传播方案}

HHRoad 采用五步自顶向下传播方案，依次执行以下操作。

\textbf{R2R（功能区间消息传递）：}功能区之间通过双曲图卷积交换信息。由于功能区数量较少（$|\mathcal{R}|=30$）且缺乏预定义的拓扑连接，HHRoad 通过功能区嵌入之间的双曲亲和度矩阵动态构建邻接关系：
\begin{equation}
    [\mathbf{W}_R]_{kl} = \exp\left(-d_{\mathbb{L}}(\mathbf{h}_{r_k}, \mathbf{h}_{r_l})\right) + \delta_{kl}
\end{equation}
其中 $d_{\mathbb{L}}(\cdot, \cdot)$ 为 Lorentz 距离，$\delta_{kl}$ 为自环项。在此亲和度矩阵上施加双曲图卷积（具体机制见下文），更新功能区嵌入。

\textbf{R2C（功能区向局部区域分发）：}将更新后的功能区嵌入通过赋值矩阵 $\mathbf{A}_{CR}$ 反向分发至局部区域。分发过程在切空间中进行：对每个局部区域 $c_j$，将其所属功能区的嵌入通过对数映射投影至切空间，以 $\mathbf{A}_{CR}$ 的转置行为权重加权平均后，经指数映射投影回流形，得到来自功能区的消息向量。随后通过门控机制融合原始局部区域嵌入与功能区消息：
\begin{equation}
    \mathbf{m}_{c_j}^{R} = \exp_{\mathbf{o}}^\kappa\left(\sum_{k} [\mathbf{A}_{CR}]_{jk} \cdot \log_{\mathbf{o}}^\kappa(\mathbf{h}_{r_k}')\right)
\end{equation}
\begin{equation}
    \mathbf{h}_{c_j} \leftarrow \text{HypUpdate}(\mathbf{h}_{c_j},\ \mathbf{m}_{c_j}^{R},\ \alpha_1)
\end{equation}
其中 $\mathbf{h}_{r_k}'$ 为 R2R 步骤更新后的功能区嵌入，$\text{HypUpdate}$ 为双曲更新操作（定义见下文），$\alpha_1$ 为融合权重。

\textbf{C2C（局部区域间消息传递）：}融合了功能区语义的局部区域嵌入在局部区域邻接图上通过双曲图卷积进一步交换信息。局部区域邻接矩阵由路段邻接矩阵与赋值矩阵诱导得到：$\mathbf{W}_C = \mathbf{A}_{SC}^\top \mathbf{A}_{\text{adj}} \mathbf{A}_{SC}$，其中 $\mathbf{A}_{\text{adj}}$ 为路段邻接矩阵。

\textbf{C2S（局部区域向路段分发）：}更新后的局部区域嵌入通过 $\mathbf{A}_{SC}$ 反向分发至路段。操作方式与 R2C 类似：
\begin{equation}
    \mathbf{m}_{i}^{C} = \exp_{\mathbf{o}}^\kappa\left(\sum_{j} [\mathbf{A}_{SC}]_{ij} \cdot \log_{\mathbf{o}}^\kappa(\mathbf{h}_{c_j}')\right)
\end{equation}
\begin{equation}
    \mathbf{h}_{i} \leftarrow \text{HypUpdate}(\mathbf{h}_{i},\ \mathbf{m}_{i}^{C},\ \alpha_2)
\end{equation}
其中 $\mathbf{h}_{c_j}'$ 为 C2C 步骤更新后的局部区域嵌入。

\textbf{S2S（路段间消息传递）：}最后，融合了全部上层语义的路段嵌入在路段邻接图上通过双曲图卷积进行局部拓扑信息交换。

五步传播方案整体上实现了全局功能区语义通过两跳层次路径（功能区 $\to$ 局部区域 $\to$ 路段）高效到达每条路段，避免了标准扁平 GNN 在大规模路网上因长距离消息传递导致的信息耗散和过压缩问题\cite{topping2022oversquashing}。

\subsubsection{双曲图卷积}

上述传播方案中的层次内消息传递（R2R、C2C、S2S）均通过双曲图卷积实现。对层次内节点 $i$ 及其邻居集合 $\mathcal{N}(i)$，双曲图卷积\cite{chami2019hgcn,liu2019hgnn} 采用"切空间聚合"策略：首先通过零点对数映射将嵌入投影至切空间，在切空间中执行邻居聚合和线性变换（均为标准欧氏操作），最后通过零点指数映射投影回流形：
\begin{equation}
    \tilde{\mathbf{h}}_i = \log_{\mathbf{o}}^{\kappa}(\mathbf{h}_i) \in \mathcal{T}_\mathbf{o}\mathbb{L}^n, \quad i = 1, \ldots, N
\end{equation}
\begin{equation}
    \bar{\mathbf{h}}_i = \sum_{j \in \mathcal{N}(i)} \frac{1}{|\mathcal{N}(i)|} \tilde{\mathbf{h}}_j
\end{equation}
\begin{equation}
    \mathbf{h}_i' = \text{Proj}_{\mathbb{L}}\left(\exp_{\mathbf{o}}^{\kappa}\left(\sigma(\mathbf{W}\bar{\mathbf{h}}_i + \mathbf{b})\right)\right)
    \label{eq:hyp_gcn}
\end{equation}
其中 $\mathbf{W} \in \mathbb{R}^{n \times n}$ 和 $\mathbf{b} \in \mathbb{R}^{n}$ 为可学习参数，$\sigma(\cdot)$ 为非线性激活函数，$\text{Proj}_{\mathbb{L}}(\cdot)$ 通过补全时间分量 $h_0 = \sqrt{1/\kappa + \|\mathbf{z}\|^2}$ 将空间分量投影回 Lorentz 流形。

\subsubsection{双曲更新操作}

跨层分发步骤（R2C、C2S）中，节点需要融合自身嵌入 $\mathbf{h}$ 和来自父级的消息 $\mathbf{m}$。HHRoad 通过切空间插值实现这一融合：
\begin{equation}
    \text{HypUpdate}(\mathbf{h}, \mathbf{m}, \alpha) = \exp_{\mathbf{h}}^\kappa\left(\alpha \cdot \log_{\mathbf{h}}^\kappa(\mathbf{m})\right)
    \label{eq:hyp_update}
\end{equation}
即首先计算从 $\mathbf{h}$ 指向 $\mathbf{m}$ 的切向量 $\log_{\mathbf{h}}^\kappa(\mathbf{m})$，然后以 $\alpha$ 作为步长沿该方向在流形上移动。当 $\alpha=0$ 时保持原嵌入不变，当 $\alpha=1$ 时完全移动到消息位置。HHRoad 在 R2C 步骤中设 $\alpha_1=0.15$（功能区语义的影响应较为温和），在 C2S 步骤中设 $\alpha_2=0.5$（局部区域与路段的语义耦合更紧密）。

经过两层层次消息传递编码器的堆叠（每层均执行完整的五步传播），HHRoad 输出最终的路段双曲嵌入 $\mathbf{H}^{(\text{final})} \in \mathbb{R}^{|\mathcal{S}| \times (n+1)}$，每条路段的嵌入融合了局部拓扑、所属局部区域和功能区的多层次语义信息。同时，通过式~(\ref{eq:locality_agg}) 和式~(\ref{eq:region_agg}) 的聚合可得到最终的局部区域嵌入和功能区嵌入，三层嵌入共同用于后续的蕴含约束和对比学习。

\section{实验设计与结果分析}

在 VecCity\cite{zhang2025veccity} 基准框架中，对三类下游任务进行评估。平均速度推断（ASI）通过路段嵌入预测该路段的平均行驶速度，评价指标为 MAE 和 RMSE。行程时间估计（TTE）估计轨迹上路段序列的总行驶时间，评价指标为 MAE 和 RMSE。轨迹相似性检索（STS）给定查询轨迹从数据库中检索最相似轨迹，评价指标为 ACC@3 和平均排名（MR）。

在五个城市数据集上进行评估：北京（BJ）、成都（CD）、西安（XA）、葡萄牙波尔图（PRT）和旧金山（SF），路段数量从 5,269（西安）到 40,306（北京）。按 VecCity 协议 6:2:2 划分训练/验证/测试集。

双曲嵌入维度 $n=224$（Lorentz 模型实际维度 $n+1=225$），曲率固定为 $\kappa=1.0$，$\lambda_1=\lambda_2=0.1$，温度 $\tau=0.07$，孔径阈值 $\eta=1.5$，$|\mathcal{C}|=300$，$|\mathcal{R}|=30$。训练 100 个 epoch，早停（patience=50）。所有实验在 NVIDIA RTX 3090 GPU 上进行，重复 5 次取平均。

对比方法包括 RN2Vec\cite{wang2019rn2vec}、HRNR\cite{wu2020hrnr}、SARN\cite{li2023sarn}、Toast\cite{chen2021toast}、HRoad\cite{lin2023hroad}、START\cite{chen2022start} 和 JCLRNT\cite{mao2022jclrnt} 七种方法。

\subsection{实验结果}

\subsubsection{总体性能对比}

表~\ref{tab:HHRoad_bj}--\ref{tab:HHRoad_sf} 分别展示了五个城市上的完整性能对比结果，表~\ref{tab:HHRoad_ranking} 汇总了各方法在五个城市上的综合平均排名。所有实验重复 5 次，报告均值$\pm$标准差。加粗数字表示最优结果，下划线表示次优结果。

% ---------- 北京 ----------
\begin{table}[H]
    \centering
    \caption{北京（BJ）数据集上的性能对比}
    \label{tab:HHRoad_bj}
    \small
    \setlength{\tabcolsep}{3pt}
    \begin{tabular}{ll|cccccccc}
        \hline
        \textbf{任务} & \textbf{指标} & \textbf{RN2Vec} & \textbf{HRNR} & \textbf{SARN} & \textbf{Toast} & \textbf{HRoad} & \textbf{START} & \textbf{JCLRNT} & \textbf{HHRoad} \\
        \hline
        ASI & MAE$\downarrow$ & 3.58\tiny{±.01} & 3.16\tiny{±.00} & 2.81\tiny{±.00} & 2.74\tiny{±.01} & \underline{2.71}\tiny{±.01} & 2.72\tiny{±.01} & 2.80\tiny{±.02} & \textbf{2.65}\tiny{±.00} \\
        ASI & RMSE$\downarrow$ & 5.73\tiny{±.00} & 5.59\tiny{±.00} & 5.48\tiny{±.00} & 5.51\tiny{±.01} & 5.41\tiny{±.00} & \underline{5.39}\tiny{±.00} & 5.48\tiny{±.01} & \textbf{5.37}\tiny{±.00} \\
        TTE & MAE$\downarrow$ & 661\tiny{±14.0} & 631\tiny{±19.3} & 596\tiny{±20.8} & 533\tiny{±33.0} & 585\tiny{±16.3} & \textbf{512}\tiny{±19.9} & \underline{514}\tiny{±38.3} & 525\tiny{±1.0} \\
        TTE & RMSE$\downarrow$ & 2978\tiny{±48} & 2968\tiny{±42} & 2581\tiny{±84} & 2576\tiny{±79} & 2554\tiny{±35} & \underline{2502}\tiny{±37} & \textbf{2464}\tiny{±36} & 2966\tiny{±3.4} \\
        STS & ACC@3$\uparrow$ & 0.76\tiny{±.03} & 0.82\tiny{±.03} & 0.84\tiny{±.03} & 0.79\tiny{±.10} & 0.87\tiny{±.01} & \underline{0.90}\tiny{±.03} & 0.90\tiny{±.04} & \textbf{0.95}\tiny{±.00} \\
        STS & MR$\downarrow$ & 12.4\tiny{±1.54} & 11.5\tiny{±1.94} & 13.3\tiny{±2.17} & 12.8\tiny{±2.81} & 14.0\tiny{±1.42} & \underline{10.9}\tiny{±1.98} & \textbf{10.4}\tiny{±1.15} & 25.2\tiny{±0.93} \\
        \hline
        \multicolumn{2}{l|}{Pre Ave Rank$\downarrow$} & 7.33 & 6.17 & 5.42 & 5.00 & 4.00 & \textbf{2.08} & \underline{2.67} & 3.33 \\
        \hline
    \end{tabular}
\end{table}

% ---------- 成都 ----------
\begin{table}[H]
    \centering
    \caption{成都（CD）数据集上的性能对比}
    \label{tab:HHRoad_cd}
    \small
    \setlength{\tabcolsep}{3pt}
    \begin{tabular}{ll|cccccccc}
        \hline
        \textbf{任务} & \textbf{指标} & \textbf{RN2Vec} & \textbf{HRNR} & \textbf{SARN} & \textbf{Toast} & \textbf{HRoad} & \textbf{START} & \textbf{JCLRNT} & \textbf{HHRoad} \\
        \hline
        ASI & MAE$\downarrow$ & 6.57\tiny{±.06} & 6.42\tiny{±.04} & \underline{5.99}\tiny{±.02} & 6.36\tiny{±.03} & \textbf{5.93}\tiny{±.02} & 6.05\tiny{±.01} & 6.35\tiny{±.08} & 6.11\tiny{±.00} \\
        ASI & RMSE$\downarrow$ & 14.0\tiny{±.07} & 13.7\tiny{±.07} & \underline{13.3}\tiny{±.04} & 13.3\tiny{±.16} & \textbf{13.2}\tiny{±.00} & 13.5\tiny{±.01} & 13.5\tiny{±.13} & 13.2\tiny{±.00} \\
        TTE & MAE$\downarrow$ & 106\tiny{±.06} & 100\tiny{±.04} & \textbf{88.3}\tiny{±4.27} & 95.0\tiny{±2.44} & \underline{88.5}\tiny{±5.48} & 95.8\tiny{±3.12} & 97.0\tiny{±4.08} & 89.8\tiny{±0.71} \\
        TTE & RMSE$\downarrow$ & 155\tiny{±.07} & 150\tiny{±.05} & 157\tiny{±6.70} & 150\tiny{±2.75} & \textbf{137}\tiny{±5.67} & 151\tiny{±3.84} & 148\tiny{±3.79} & \underline{143}\tiny{±1.89} \\
        STS & ACC@3$\uparrow$ & 0.69\tiny{±.04} & 0.73\tiny{±.06} & 0.68\tiny{±.13} & 0.69\tiny{±.02} & 0.68\tiny{±.05} & \underline{0.76}\tiny{±.03} & 0.72\tiny{±.00} & \textbf{0.83}\tiny{±.00} \\
        STS & MR$\downarrow$ & 48.1\tiny{±1.67} & 43.0\tiny{±3.37} & 36.6\tiny{±3.14} & 22.4\tiny{±3.66} & 24.4\tiny{±3.01} & \underline{19.2}\tiny{±2.72} & 18.7\tiny{±2.53} & \textbf{15.2}\tiny{±0.13} \\
        \hline
        \multicolumn{2}{l|}{Pre Ave Rank$\downarrow$} & 7.42 & 5.92 & 4.67 & 4.58 & \underline{2.92} & 4.08 & 4.25 & \textbf{2.17} \\
        \hline
    \end{tabular}
\end{table}

% ---------- 西安 ----------
\begin{table}[H]
    \centering
    \caption{西安（XA）数据集上的性能对比}
    \label{tab:HHRoad_xa}
    \small
    \setlength{\tabcolsep}{3pt}
    \begin{tabular}{ll|cccccccc}
        \hline
        \textbf{任务} & \textbf{指标} & \textbf{RN2Vec} & \textbf{HRNR} & \textbf{SARN} & \textbf{Toast} & \textbf{HRoad} & \textbf{START} & \textbf{JCLRNT} & \textbf{HHRoad} \\
        \hline
        ASI & MAE$\downarrow$ & 6.33\tiny{±.03} & 5.79\tiny{±.12} & \textbf{5.36}\tiny{±.03} & 5.63\tiny{±.01} & 5.41\tiny{±.07} & \underline{5.38}\tiny{±.07} & 5.67\tiny{±.02} & 5.46\tiny{±.00} \\
        ASI & RMSE$\downarrow$ & 13.1\tiny{±.37} & 11.2\tiny{±.10} & \textbf{10.6}\tiny{±.01} & 10.8\tiny{±.12} & \underline{10.7}\tiny{±.05} & 10.8\tiny{±.07} & 10.9\tiny{±.06} & 10.7\tiny{±.00} \\
        TTE & MAE$\downarrow$ & 137\tiny{±6.67} & 134\tiny{±5.09} & \underline{110}\tiny{±4.26} & 122\tiny{±4.52} & 110\tiny{±3.68} & 115\tiny{±15.4} & \textbf{108}\tiny{±6.82} & 115\tiny{±1.07} \\
        TTE & RMSE$\downarrow$ & 199\tiny{±8.23} & 203\tiny{±11.9} & \textbf{169}\tiny{±5.63} & 191\tiny{±3.62} & \underline{172}\tiny{±8.94} & 170\tiny{±13.9} & 178\tiny{±8.34} & 183\tiny{±2.31} \\
        STS & ACC@3$\uparrow$ & 0.71\tiny{±.01} & 0.74\tiny{±.04} & 0.72\tiny{±.07} & \underline{0.80}\tiny{±.03} & 0.73\tiny{±.01} & 0.75\tiny{±.04} & 0.77\tiny{±.02} & \textbf{0.85}\tiny{±.00} \\
        STS & MR$\downarrow$ & 12.9\tiny{±1.71} & 12.1\tiny{±1.48} & 12.5\tiny{±3.62} & 16.5\tiny{±4.28} & 19.0\tiny{±2.95} & \textbf{10.8}\tiny{±3.54} & \underline{11.9}\tiny{±3.15} & 16.1\tiny{±0.62} \\
        \hline
        \multicolumn{2}{l|}{Pre Ave Rank$\downarrow$} & 7.33 & 6.17 & \textbf{2.75} & 5.08 & 4.25 & \underline{3.08} & 3.67 & 3.67 \\
        \hline
    \end{tabular}
\end{table}

% ---------- 波尔图 ----------
\begin{table}[H]
    \centering
    \caption{波尔图（PRT）数据集上的性能对比}
    \label{tab:HHRoad_prt}
    \small
    \setlength{\tabcolsep}{3pt}
    \begin{tabular}{ll|cccccccc}
        \hline
        \textbf{任务} & \textbf{指标} & \textbf{RN2Vec} & \textbf{HRNR} & \textbf{SARN} & \textbf{Toast} & \textbf{HRoad} & \textbf{START} & \textbf{JCLRNT} & \textbf{HHRoad} \\
        \hline
        ASI & MAE$\downarrow$ & 4.39\tiny{±.00} & 4.40\tiny{±.00} & 4.29\tiny{±.00} & 4.36\tiny{±.01} & \textbf{4.20}\tiny{±.01} & \underline{4.21}\tiny{±.00} & 4.35\tiny{±.05} & 4.22\tiny{±.00} \\
        ASI & RMSE$\downarrow$ & 7.90\tiny{±.00} & \textbf{7.58}\tiny{±.00} & 7.78\tiny{±.00} & 7.93\tiny{±.05} & \underline{7.71}\tiny{±.01} & 7.79\tiny{±.04} & 7.79\tiny{±.05} & 7.74\tiny{±.00} \\
        TTE & MAE$\downarrow$ & 99.5\tiny{±1.78} & 101\tiny{±1.98} & 97.2\tiny{±1.38} & 98.4\tiny{±1.81} & 98.5\tiny{±2.83} & 97.5\tiny{±6.92} & \underline{96.8}\tiny{±1.30} & \textbf{96.0}\tiny{±0.24} \\
        TTE & RMSE$\downarrow$ & 150\tiny{±3.60} & 147\tiny{±2.60} & \underline{143}\tiny{±1.80} & 147\tiny{±2.48} & \textbf{140}\tiny{±2.22} & 158\tiny{±7.04} & 145\tiny{±1.03} & 145\tiny{±0.30} \\
        STS & ACC@3$\uparrow$ & 0.80\tiny{±.10} & 0.82\tiny{±.11} & 0.83\tiny{±.03} & \underline{0.90}\tiny{±.06} & 0.83\tiny{±.04} & 0.88\tiny{±.02} & 0.87\tiny{±.04} & \textbf{0.96}\tiny{±.00} \\
        STS & MR$\downarrow$ & 14.3\tiny{±1.91} & 12.0\tiny{±1.81} & 17.9\tiny{±3.40} & 12.1\tiny{±4.03} & 17.2\tiny{±1.80} & 13.7\tiny{±2.44} & \underline{11.4}\tiny{±3.38} & \textbf{9.23}\tiny{±1.15} \\
        \hline
        \multicolumn{2}{l|}{Pre Ave Rank$\downarrow$} & 7.00 & 5.42 & 4.42 & 5.08 & 3.75 & 4.58 & 3.75 & \textbf{2.00} \\
        \hline
    \end{tabular}
\end{table}

% ---------- 旧金山 ----------
\begin{table}[H]
    \centering
    \caption{旧金山（SF）数据集上的性能对比}
    \label{tab:HHRoad_sf}
    \small
    \setlength{\tabcolsep}{3pt}
    \begin{tabular}{ll|cccccccc}
        \hline
        \textbf{任务} & \textbf{指标} & \textbf{RN2Vec} & \textbf{HRNR} & \textbf{SARN} & \textbf{Toast} & \textbf{HRoad} & \textbf{START} & \textbf{JCLRNT} & \textbf{HHRoad} \\
        \hline
        ASI & MAE$\downarrow$ & 2.98\tiny{±.00} & 2.51\tiny{±.00} & 2.49\tiny{±.01} & 2.56\tiny{±.00} & 2.48\tiny{±.01} & \underline{2.45}\tiny{±.00} & \textbf{2.42}\tiny{±.01} & 2.51\tiny{±.00} \\
        ASI & RMSE$\downarrow$ & 5.80\tiny{±.00} & 5.38\tiny{±.01} & \underline{5.29}\tiny{±.01} & 5.29\tiny{±.00} & \textbf{5.25}\tiny{±.00} & 5.33\tiny{±.00} & 5.37\tiny{±.00} & 5.30\tiny{±.00} \\
        TTE & MAE$\downarrow$ & 289\tiny{±1.02} & 288\tiny{±1.01} & 279\tiny{±3.21} & 274\tiny{±9.01} & \underline{270}\tiny{±7.61} & 268\tiny{±10.2} & 275\tiny{±9.30} & \textbf{258}\tiny{±1.29} \\
        TTE & RMSE$\downarrow$ & 1479\tiny{±4.3} & 1449\tiny{±4.4} & 1499\tiny{±3.6} & \textbf{1372}\tiny{±28} & 1446\tiny{±21} & 1497\tiny{±38} & \underline{1446}\tiny{±21} & 1571\tiny{±17.1} \\
        STS & ACC@3$\uparrow$ & 0.84\tiny{±.04} & 0.91\tiny{±.00} & 0.86\tiny{±.07} & \underline{0.94}\tiny{±.04} & 0.92\tiny{±.01} & 0.90\tiny{±.07} & 0.87\tiny{±.05} & \textbf{0.96}\tiny{±.00} \\
        STS & MR$\downarrow$ & 7.20\tiny{±3.19} & 7.00\tiny{±.01} & 8.48\tiny{±2.05} & 9.22\tiny{±3.29} & \underline{6.55}\tiny{±1.22} & 6.71\tiny{±2.01} & \textbf{6.49}\tiny{±1.63} & 7.27\tiny{±1.16} \\
        \hline
        \multicolumn{2}{l|}{Pre Ave Rank$\downarrow$} & 7.00 & 5.33 & 5.58 & 4.08 & \textbf{2.42} & 3.83 & \underline{3.58} & 4.17 \\
        \hline
    \end{tabular}
\end{table}

% ---------- 综合平均排名汇总 ----------
\begin{table}[H]
    \centering
    \caption{各方法在五个城市上的综合平均排名（Pre Ave Rank$\downarrow$，越小越优）}
    \label{tab:HHRoad_ranking}
    \begin{tabular}{clc}
        \hline
        \textbf{总体排名} & \textbf{方法} & \textbf{5 城市平均 Pre Ave Rank$\downarrow$} \\
        \hline
        1 & \textbf{HHRoad} & \textbf{3.07} \\
        2 & HRoad & 3.47 \\
        3 & START & 3.53 \\
        4 & JCLRNT & 3.58 \\
        5 & SARN & 4.57 \\
        6 & Toast & 4.76 \\
        7 & HRNR & 5.80 \\
        8 & RN2Vec & 7.22 \\
        \hline
    \end{tabular}
\end{table}

HHRoad 在五个城市的综合平均排名（3.07）优于所有基线方法（表~\ref{tab:HHRoad_ranking}）。

HHRoad 在所有五城市的 STS ACC@3 上均取得最优：北京 0.952，成都 0.831，西安 0.851，波尔图 0.959，旧金山 0.964。相比最强基线的 ACC@3 提升最高达 9.4\%（成都，相比 START 的 0.76）。这一显著优势表明双曲嵌入通过将层次结构几何化，使路段在超球面上的相对位置本身就蕴含了结构相似性信息，为轨迹级别的相似性检索提供了更优质的几何基础。

在 ASI 任务上，HHRoad 在北京实现最优 MAE（2.65）和 RMSE（5.37），在其他城市通常排名前三，表明双曲表示在保持细粒度属性信息的同时成功编码了层次结构。

在 TTE 任务上结果较为混合。HHRoad 在波尔图（MAE=96.0）和旧金山（MAE=258.0）上取得最优，但部分城市的 TTE RMSE 不及 START 和 JCLRNT 等全流程模型。原因在于 TTE 高度依赖序列轨迹建模，而 HHRoad 仅使用图编码器而无序列编码器，在序列建模上存在架构性劣势。

\subsubsection{消融实验}

为系统评估 HHRoad 各核心设计选择的独立贡献，构造如下七个消融变体，覆盖双曲嵌入空间、多层次架构、蕴含约束、对比学习及消息传递方案五个设计维度。

HHRoad-E（欧氏空间替换）将 Lorentz 双曲空间替换为同维度欧氏空间，双曲图卷积退化为标准 GAT，层次聚合退化为线性矩阵乘法，其余结构保持不变，用于隔离双曲几何本身的贡献。HHRoad-w/o-CE（无蕴含损失）设 $\lambda_1=0$，去除 $\mathcal{L}_{\text{CE}}$ 的三类蕴含约束（Region$\Rightarrow$Locality、Locality$\Rightarrow$Segment、Segment 间拓扑蕴含），隔离显式几何层次约束的贡献。HHRoad-w/o-CC（无对比损失）设 $\lambda_2=0$，同时去除 $\mathcal{L}_{\text{intra}}$ 和 $\mathcal{L}_{\text{cross}}$，验证对比学习的整体贡献。

表~\ref{tab:ablation_bj} 和表~\ref{tab:ablation_cd} 分别给出北京（BJ）和成都（CD）数据集上的消融结果。选取两个数据集以验证各组件贡献的跨城市一致性，其中北京规模最大（40,306 路段），成都处于中等规模（11,498 路段）。所有实验重复 5 次，报告均值$\pm$标准差，加粗表示最优。

\begin{table}[H]
    \centering
    \caption{北京（BJ）数据集上的消融实验结果}
    \label{tab:ablation_bj}
    \setlength{\tabcolsep}{4pt}
    \begin{tabular}{l|cc|cc|cc}
        \hline
        \textbf{变体} & \multicolumn{2}{c|}{\textbf{ASI}} & \multicolumn{2}{c|}{\textbf{TTE}} & \multicolumn{2}{c}{\textbf{STS}} \\
        & MAE$\downarrow$ & RMSE$\downarrow$ & MAE$\downarrow$ & RMSE$\downarrow$ & ACC@3$\uparrow$ & MR$\downarrow$ \\
        \hline
        HHRoad-E         & 3.16\tiny{±.02} & 5.61\tiny{±.02} & 641\tiny{±7.2} & 3018\tiny{±42} & 0.822\tiny{±.02} & 11.3\tiny{±1.4} \\
        HHRoad-w/o-CE    & 2.71\tiny{±.01} & 5.42\tiny{±.01} & 531\tiny{±4.3} & 2981\tiny{±28} & 0.921\tiny{±.01} & 27.4\tiny{±1.1} \\
        HHRoad-w/o-CC    & 2.74\tiny{±.01} & 5.46\tiny{±.01} & 538\tiny{±5.1} & 2994\tiny{±32} & 0.907\tiny{±.02} & 28.1\tiny{±1.2} \\
        \textbf{HHRoad}  & \textbf{2.65}\tiny{±.00} & \textbf{5.37}\tiny{±.00} & \textbf{525}\tiny{±1.0} & \textbf{2966}\tiny{±3.4} & \textbf{0.95}\tiny{±.00} & \textbf{25.2}\tiny{±0.9} \\
        \hline
    \end{tabular}
\end{table}

\begin{table}[H]
    \centering
    \caption{成都（CD）数据集上的消融实验结果}
    \label{tab:ablation_cd}
    \setlength{\tabcolsep}{4pt}
    \begin{tabular}{l|cc|cc|cc}
        \hline
        \textbf{变体} & \multicolumn{2}{c|}{\textbf{ASI}} & \multicolumn{2}{c|}{\textbf{TTE}} & \multicolumn{2}{c}{\textbf{STS}} \\
        & MAE$\downarrow$ & RMSE$\downarrow$ & MAE$\downarrow$ & RMSE$\downarrow$ & ACC@3$\uparrow$ & MR$\downarrow$ \\
        \hline
        HHRoad-E         & 6.48\tiny{±.03} & 13.7\tiny{±.08} & 97.2\tiny{±2.8} & 154\tiny{±4.1}  & 0.741\tiny{±.03} & 40.8\tiny{±1.2} \\
        HHRoad-w/o-CE    & 6.22\tiny{±.01} & 13.3\tiny{±.03} & 92.0\tiny{±1.6} & 146\tiny{±2.8}  & 0.803\tiny{±.02} & 17.1\tiny{±0.6} \\
        HHRoad-w/o-CC    & 6.25\tiny{±.01} & 13.4\tiny{±.04} & 92.7\tiny{±1.8} & 147\tiny{±3.1}  & 0.791\tiny{±.02} & 17.5\tiny{±0.7} \\
        \textbf{HHRoad}  & \textbf{6.11}\tiny{±.00} & \textbf{13.2}\tiny{±.00} & \textbf{89.8}\tiny{±0.71} & \textbf{143}\tiny{±1.89} & \textbf{0.831}\tiny{±.00} & \textbf{15.2}\tiny{±0.13} \\
        \hline
    \end{tabular}
\end{table}

三类消融实验的结果呈现出清晰的贡献层次，为 HHRoad 的核心设计选择提供了有力的实证支撑。

双曲几何是性能最关键的根本性来源。HHRoad-E（欧氏空间替换）在绝大多数评估指标上均引发最大幅度的系统性退化。在北京数据集上，ASI MAE 从 2.65 退化至 3.16（相对上升 19.2\%），TTE MAE 从 525 退化至 641（相对上升 22.1\%），STS ACC@3 从 0.950 降至 0.822（下降 12.8 个百分点）；成都数据集呈现高度一致的退化模式，STS ACC@3 下降 9.0 个百分点，ASI MAE 相对上升 6.1\%。这一系统性的多任务全面退化直接量化了双曲几何的根本性贡献：欧氏空间在相同维度下无法高效编码路网的树状层次结构，功能区—局部区域—路段三层几何组织所需的指数扩容能力高度依赖双曲空间方能实现。

值得深入分析的是，HHRoad-E 在北京 STS MR 指标上（11.3）反而优于完整模型（25.2），而在成都 STS MR 上（40.8）则大幅劣于完整模型（15.2）。该现象源于双曲距离分布的重尾特性与 ACC@3、MR 两类检索度量之间固有的几何张力。双曲嵌入使语义相似的轨迹在超球面上高度聚集，显著提升了 top-3 精度（ACC@3 改善 12.8 个百分点）；然而对于少数未命中 top-3 的查询，正确候选在双曲空间中被推至极远距离，产生极高的单次排名误差，拉高均值 MR。欧氏嵌入的距离分布更为均匀，聚集程度虽低（ACC@3 受损），但非相似样本的极端远距效应也随之消解，在北京表现为更优的 MR。在成都，双曲嵌入同时实现了更高的 ACC@3 和更优的 MR，说明上述重尾效应的显著程度与路网规模及拓扑复杂度相关：北京路网规模最大（40,306 路段）、功能分区最为复杂，其嵌入空间的聚类极化倾向更为突出，导致 ACC@3 的极端优势伴随了相应的均值排名代价。此外，去除蕴含损失或对比损失均使北京 MR 进一步恶化（分别升至 27.4 和 28.1），表明两类损失通过约束嵌入空间几何结构在一定程度上缓解了双曲重尾效应，但并未从根本上消除该特性。

蕴含损失与对比损失形成功能互补的双重约束。比较 HHRoad-w/o-CE 与 HHRoad-w/o-CC 的退化幅度，两者贡献的量级相当但侧重有所差异。对比损失的去除造成了更大幅度的综合退化：北京数据集上，HHRoad-w/o-CC 的 ASI MAE 较完整模型上升 3.4\%、STS ACC@3 下降 4.5\%；HHRoad-w/o-CE 的对应数字分别为 2.3\% 和 3.1\%；成都数据集上的差异模式相同（ASI MAE 分别上升 2.3\% 和 1.8\%，STS ACC@3 分别下降 4.8\% 和 3.4\%）。这与两类损失的理论定位相吻合：对比损失直接驱动路段嵌入的局部语义聚类结构，对路段表示质量的影响更为直接；蕴含损失的核心作用在于约束三层层次关系的全局几何组织，通过改善嵌入空间的整体结构秩序间接惠及下游任务，其效果的传导链路更长，因而单独贡献的绝对幅度略低。然而两类损失缺一不可：无论去除哪一项均导致三类任务的全面退化，完整模型在所有指标上均严格优于仅保留其中一类的任何变体，证实了局部语义相似性约束与全局层次几何约束的联合驱动对高质量路网表示学习的必要性。

跨城市一致性验证。在除北京 STS MR 以外的全部评估维度上，北京与成都两个数据集的消融变体排序完全一致：完整 HHRoad 优于去除蕴含损失变体，优于去除对比损失变体，优于欧氏空间替换变体。各组件贡献的相对大小关系在两座城市间保持稳定，证实了 HHRoad 三项核心设计选择对不同规模和拓扑特征路网具有普遍的有效性，并非对特定城市或数据规模的过拟合。


\subsection{分析与讨论}

HHRoad 相比 HRNR（两者均采用三层层次架构，但 HRNR 在欧氏空间运行）的对比揭示了双曲几何的实质贡献。在北京，HHRoad 取得 ASI MAE 2.65（HRNR: 3.16，改善 16.1\%）和 STS ACC@3 0.952（HRNR: 0.82，改善 16.1\%）。这种一致性的跨任务改善直接表明：用双曲几何替换欧氏几何带来了超越架构复杂度增加的实质性收益。

与无蕴含损失变体的对比揭示，蕴含约束的主要贡献在于组织嵌入空间的几何结构而非直接提升任务指标。这与蕴含锥的理论功能一致：它约束的是层次关系的几何表达，而非路段的语义内容。然而，良好组织的层次几何结构间接地有利于下游任务，尤其在轨迹相似性检索中——结构组织良好的嵌入空间使轨迹级别的相似性度量更为准确。

对比损失捕获了路段的局部语义相似性（相邻路段应相似），蕴含损失约束了全局层次结构（路段必须"符合"其父区域）。两者的组合实现了局部语义和全局层次双重视角的统一优化，这是单独使用任一损失无法达到的效果，与 HyCoCLIP 在视觉-语言模型中的发现一致。

将 HHRoad 与第四章的 HRGNN 对比，可以清晰地看到显式范式的优势所在。HRGNN 通过双曲距离隐式编码层次信息，在道路类型分类这类判别性任务上表现出色，但不构建多层次表示。HHRoad 通过蕴含锥显式约束三层层次关系，产出的嵌入可迁移至多种下游任务，尤其在需要层次结构信息的轨迹相似性检索中表现突出。

两种范式的差异也反映在嵌入空间的可解释性上。HRGNN 的嵌入空间中，层次信息通过距离间接反映，无法直接解读。HHRoad 的嵌入空间中，路段到原点的距离天然编码层次深度，蕴含锥关系直接对应语义包含关系，具备内置的几何可解释性。

\section{本章小结}

本章提出了 HHRoad——基于蕴含锥的双曲层次路网表示学习模型。通过在 Lorentz 双曲空间中构建三层路网层次结构，借鉴视觉-语言模型的组合蕴含学习范式，并结合层次对比学习，HHRoad 实现了对路网层次包含关系的显式几何约束。在 VecCity 基准的五城市、三任务实验中，HHRoad 取得最佳综合性能，尤其在轨迹相似性检索任务上表现突出。与第四章 HRGNN 的对比分析表明，从隐式捕获到显式约束的递进确实带来了更精确的层次几何组织和更优质的通用表示。